\chapter{Pain}

\section*{Definition}
Unpleasant sensory and emotional experience associated with actual or potential tissue damage.

\section*{Pathophysiology}
Diverse etiations across all body systems; characterization essential for diagnosis including location, quality, radiation, severity, timing.

\section*{Etiology}
\begin{itemize}
  \item Various depending on location and etiology
\end{itemize}

\section*{Indications for Evaluation}
Seek care if:
\begin{itemize}
  \item Severe or worsening symptoms
  \item Symptoms lasting more than Persistent/severe or interfering with function requires medical evaluation
  \item Accompanied by other concerning signs
\end{itemize}

\section*{Related Symptoms}
\begin{itemize}
  \item Varies greatly based on type and location
\end{itemize}
\textbf{Note:} Always consider serious underlying conditions when evaluating persistent pain.

\section*{Self-Care / Home Management}
\begin{itemize}
  \item Rest the affected area and avoid activities that aggravate the pain.
  \item Apply ice packs for 15-20 minutes at a time during the first 48 hours, then switch to heat if preferred.
  \item Over-the-counter analgesics such as acetaminophen or ibuprofen may provide relief (follow label instructions).
  \item Gentle stretching or movement within pain tolerance helps prevent stiffness.
  \item Seek medical evaluation if pain persists beyond 7 days despite self-care measures.
\end{itemize}

\section*{Diagnostic Approach}
\begin{itemize}
  \item A focused history assesses onset, location, quality, severity, duration, radiation, aggravating/relieving factors, and prior treatments.
  \item Physical examination includes inspection for swelling/deformity, palpation for tenderness point, range-of-motion testing, and neurovascular assessment.
  \item Imaging (X-ray, ultrasound, MRI, CT) may be indicated based on suspected etiology and red flag findings.
  \item Laboratory studies (CBC, CRP, ESR) help evaluate for inflammatory or infectious causes.
\end{itemize}

\section*{Treatment / Management Overview}
\begin{itemize}
  \item First-line management includes activity modification, physical therapy, and analgesics (NSAIDs or acetaminophen).
  \item Specific diagnoses may require targeted therapy: antibiotics for infection, corticosteroids for inflammation, surgery for structural pathology.
  \item Referral to a specialist (orthopedics, rheumatology, pain medicine) is indicated for refractory cases or complex diagnoses.
  \item Multimodal approaches combining medication, manual therapy, exercise, and behavioral strategies yield best outcomes for chronic pain.
\end{itemize}

\clearpage
